\subsubsection{イベント駆動設計}

\term{イベント駆動設計}{Event-Driven Design} は、イベントに反応して処理を行う設計方法である。イベントとは、利用者操作、メッセージ受信、状態変化、タイマー発火、センサー入力などである。

イベント駆動設計では、発行・購読方式がよく使われる。送信側はイベントを発行し、受信側は関心のあるイベントを購読する。これにより、送信側と受信側の結合を弱められる。大規模で拡張性が必要なシステムでは有効だが、イベントの順序、重複、遅延、失敗時の影響を慎重に設計する必要がある。
